\begin{abstract}
Forecasting requires calibrated probabilities grounded in how evidence maps to
outcomes: the same observation implies different futures under different
relationships. We ask whether language models (LMs) recover those relationships
or rely on surface regularities. Across real-world and controlled settings,
models distinguish probative evidence from topical controls, use context to
choose among competing relationships, and revise those choices when direct
observations arrive. Episode-specific conventions and controlled transfer extend this pattern
beyond familiar cues and settings. Outcome-based
training improves held-out market forecasts, and several checkpoints reach the
crowd-level floor. These results establish
context-dependent inductive reasoning as an observable forecasting behavior
within the tested tasks. Evaluation should measure aggregate accuracy and whether
forecasts change for the right structural reason.
\end{abstract}
